\documentclass[a4,11pt]{aleph-notas}

% -- Paquetes adicionales
\usepackage{enumitem}
\usepackage{aleph-comandos}

% -- Datos
\institucion{Facultad de Ciencias Exactas, Naturales y Ambientales}
\carrera{Catálogo STEM}
\asignatura{Álgebra lineal}
\tema{Resumen 03: Matriz inversa}
\autor[A. Merino]{Andrés Merino}
\fecha{Periodo 2026-1}

\logouno[0.16\textwidth]{Logos/logoPUCE_04_ac}
\definecolor{colortext}{HTML}{1957d1}
\definecolor{colordef}{HTML}{1957d1}
\fuente{montserrat}

\begin{document}

\encabezado

%%%%%%%%%%%%%%%%%%%%%%%%%%%%%%%%%%%%%%
\section{Inversa de una matriz}
%%%%%%%%%%%%%%%%%%%%%%%%%%%%%%%%%%%%%%

\begin{defi}
    Sea $A\in \Mat[\K]{n}{n}$ una matriz no singular o invertible si existe una matriz $B\in \Mat[\K]{n}{n}$ tal que 
    \[
        AB = BA = I_n.
    \]
    A la matriz $B$ se la denomina inversa de $A$ y se la denota por $A^{-1}$. Si no existe tal matriz, entonces se dice que $A$ es singular o no invertible. 
\end{defi}

\begin{teo}
    Si una matriz tiene inversa, la inversa es única. 
\end{teo}

\begin{teo}
    Sea $A\in \Mat[\K]{n}{n}$. 
    \begin{itemize}
    \item 
        Si $A$ es una matriz no singular, entonces $A^{-1}$ es no singular y 
        \[ 
            (A^{-1})^{-1} = A.
        \]
    \item 
        Si $A$ y $B$ son matrices no singulares, entonces $AB$ es no singular y 
        \[ 
            (AB)^{-1} = B^{-1}A^{-1}.
        \]
    \item 
        Si $A$ es una matriz no singular, entonces
        \[
            {(A^\intercal)}^{-1} = {(A^{-1})}^\intercal.
        \]
    \end{itemize}
\end{teo}

\begin{teo}
    Sean $p\in\N^*$ y $A_1, A_2, \ldots, A_p \in \Mat[\K]{n}{n}$ matrices no singulares. Se tiene que $A_1 A_2 \cdots A_p$ es no singular y 
    \[
        (A_1 A_2 \cdots A_p)^{-1} = A_p^{-1} A_{p-1}^{-1} \cdots A_1^{-1}.
    \]
\end{teo}

\begin{teo}
    Sean $A, B \in \Mat[\K]{n}{n}$. Se tiene que si $AB=I_n$, entonces $BA=I_n$.
\end{teo}

\begin{teo}
    Sea $A \in \Mat[\K]{n}{n}$, se tiene que $A$ es no singular si y solo si es equivalente por filas a $I_n$. Es más
    \[
        (A|I_n) \sim (I_n|A^{-1}).
    \]
\end{teo}

\begin{teo}
    Sea $A \in \Mat[\K]{n}{n}$, el sistema homogéneo 
    \[
        Ax = 0
    \]
    tiene una solución no trivial si y solo si $A$ es singular. 
\end{teo}

\begin{teo}
    Sea $A \in \Mat[\K]{n}{n}$. Se tiene que $A$ es no singular si y solo si el sistema lineal $Ax=b$ tiene una solución única para cada vector $b \in \K^{n}$.
\end{teo}

\begin{teo}
    Sea $A \in \Mat[\K]{n}{n}$, se tienen que las siguientes son equivalentes:
    \begin{enumerate}
    \item 
        $A$ es no singular;
    \item 
        el sistema $Ax=0$ solamente tiene la solución trivial;
    \item 
        $A$ es equivalente por filas a $I_n$;
    \item 
        $\rang(A)=n$; y 
    \item 
        el sistema lineal $Ax=b$ tiene una solución única para cada vector $b \in \K^{n}$. 
    \end{enumerate}
\end{teo}

\section{Aplicaciones de sistemas de ecuaciones}

%%%%%%%%%%%%%%%%%%%%%%%%%%%%%%%%%%%%%%
\subsection{Factorización LU}
%%%%%%%%%%%%%%%%%%%%%%%%%%%%%%%%%%%%%%

Sea $U\in\Mat[\K]{n}{n}$ una matriz triangular superior cuyas entradas en la diagonal sean distintas de $0$, dado $b\in\K^n$, el sistema 
\[
    Ux = b,
\]
cuya matriz ampliada es
\[
    \begin{pmatrix}
        u_{11} & u_{12} & u_{13} & \cdots & u_{1n} & | & b_1 \\
        0      & u_{22} & u_{23} & \cdots & u_{2n} & | & b_2 \\
        0      & 0      & u_{33} & \cdots & u_{3n} & | & b_3 \\
        \vdots & \vdots & \vdots & \ddots & \vdots & | & \vdots \\
        0      & 0      & 0      & \cdots & u_{nn} & | & b_n 
    \end{pmatrix},
\]
se puede resolver por \emph{sustitución regresiva}, obteniendo que la solución del sistema es
\[
    x_n = \frac{b_n}{u_{nn}}
    \texty
    x_k = \dfrac{b_k - \displaystyle\sum_{j=k+1}^n u_{kj}x_j}{u_{kk}},
\]
para $k=n-1,\ldots,1$.

De manera análoga, sea $L\in\Mat[\K]{n}{n}$ una matriz triangular inferior cuyas entradas en la diagonal sean distintas de $0$, dado $b\in\K^n$, el sistema 
\[
    Lx = b,
\]
cuya matriz ampliada es
\[
    \begin{pmatrix}
        l_{11} & 0      & 0      & \cdots & 0      & | & b_1 \\
        l_{21} & l_{22} & 0      & \cdots & 0      & | & b_2 \\
        l_{31} & l_{32} & l_{33} & \cdots & 0      & | & b_3 \\
        \vdots & \vdots & \vdots & \ddots & \vdots & | & \vdots \\
        l_{n1} & l_{n2} & l_{n3} & \cdots & l_{nn} & | & b_n 
    \end{pmatrix},
\]
se puede resolver por \emph{sustitución progresiva}, obteniendo que la solución del sistema es
\[
    x_1 = \frac{b_1}{l_{11}}
    \texty
    x_k = \dfrac{b_k - \displaystyle\sum_{j=1}^{k-1} l_{kj}x_j}{l_{kk}},
\]
para $k=2,\ldots,n$.

Finalmente, si $A \in \Mat[\K]{n}{n}$ se puede escribir como el producto de una matriz triangular inferior $L$ y una matriz triangular superior $U$, es decir, 
\[
    A = LU,
\]
entonces se dice que $A$ tiene una factorización $LU$, y dado $b\in\K^n$, se tiene que el sistema
\[
    Ax=b
\]
equivale al sistema
\[
    L(Ux) = b,
\]
con esto, podríamos resolver el sistema inicial, resolviendo los sistemas
\[
    Ly = b
    \texty
    Ux = y.
\]


\begin{teo}[Factorización LU]
    Sean $A\in \Mat[\K]{n}{n}$ una matriz tal que sea equivalente por filas a una matriz $U\in \Mat[\K]{n}{n}$ triangular superior, entonces se tiene que
    \[
        U = E_{m}\cdots E_2 E_1 A,
    \]
    donde $E_k$, para $k=1,\ldots,m$, son las matrices elementales correspondientes a las operaciones por filas. Si no se ha utilizado intercambio de filas y únicamente se han eliminado los elementos bajo la diagonal, entonces
    \[
        L = E_{1}^{-1} E_2^{-1} \cdots E_m^{-1}
    \]
    es una matriz triangular inferior y se tiene que $A=LU$.
\end{teo}


%%%%%%%%%%%%%%%%%%%%%%%%%%%%%%%%%%%%%%
\subsection{Mínimos cuadrados}
%%%%%%%%%%%%%%%%%%%%%%%%%%%%%%%%%%%%%%

\begin{defi}[Problema de la recta de mínimos cuadrados]
    En $\R^2$, dados $n$ puntos $(x_1,y_1)$, $(x_2,y_2)$, \ldots, $(x_n,y_n)$, se plantea encontrar la ecuación de la recta que minimice la suma de los cuadrados de las distancias verticales de los puntos a la recta, es decir, se busca la pendiente $m$ y el intercepto $b$ de una recta que minimice la función definida por
    \[
        L(b,m)=\sum_{k=1}^n(y_k-mx_k-b)^2.
    \]
    A la recta de pendiente $m$ e intercepto $b$ que minimizan la función anterior se la llama recta de mínimos cuadrados.
\end{defi}

\begin{teo}
    Dados los puntos $(x_1,y_1)$, $(x_2,y_2)$, \ldots, $(x_n,y_n)$, definimos
    \[
        y = \begin{pmatrix}
            y_1 \\ y_2 \\ \vdots \\ y_n
        \end{pmatrix},
        \qquad
        A = \begin{pmatrix}
            1 & x_1 \\ 1 & x_2 \\ \vdots & \vdots \\ 1 & x_n
        \end{pmatrix}
        \texty
        u = \begin{pmatrix}
            b \\ m
        \end{pmatrix}.
    \]
    Se tiene que el problema de la recta de mínimos cuadrados es equivalente a resolver el sistema
    \[
        A^\intercal A u = A^\intercal y.
    \]
\end{teo}

\begin{teo}
    Sea $A\in\Mat[\K]{m}{n}$ tal que $\rang(A)=n$, entonces, $A^\intercal A$ es no singular.
\end{teo}

\end{document}